\documentclass[aspectratio=169]{beamer}
\usetheme{metropolis}
\metroset{sectionpage=none, numbering=counter, progressbar=frametitle}
\usepackage{amsmath, amssymb}
\usepackage[T1]{fontenc}

\title{Domain-Partition Algorithm for Turbomachinery Blades}
\subtitle{Slide 7: Streamline Integration --- Runge-Kutta-Heun}
\date{\today}
\author{}
\institute{}

\begin{document}

\begin{frame}{Streamline Integration --- Runge-Kutta-Heun}

\begin{center}
\Large
\[
\mathbf{p}_{n+1} = \mathbf{p}_n + \frac{1}{2}(\mathbf{k}_1 + \mathbf{k}_2)
\]
\end{center}

\vspace{0.2em}
\begin{center}
\normalsize
\(
\mathbf{k}_1 = h\,\mathbf{v}(\mathbf{p}_n), \qquad
\mathbf{k}_2 = h\,\mathbf{v}(\mathbf{p}_n + \mathbf{k}_1)
\)
\end{center}

\vspace{0.2em}
\begin{itemize}
    \item \textbf{Integrate} along representative vectors emanating from singularities
    \item \textbf{RK-Heun:} 2nd-order accurate, stable near singularities
    \item \textbf{Terminate} on boundary, singularity, or $c_0$ corner
\end{itemize}

\vspace{0.3em}
\small
\textit{Code:} \texttt{StreamlineGenerator\_v2.runge\_kutta\_heun\_integrate\_streamline} \quad
\textit{Lit:} Kowalski 2015, ``Blocking piecewise-linear domains'' (Prop.~9)

\note{
Speaker notes:
Streamline integration traces curves starting from each singularity along the frame field. We use the Runge-Kutta-Heun method, a two-stage predictor-corrector scheme. k1 is the Euler step, k2 evaluates the vector field at the predicted endpoint, and the final update averages the two slopes. This gives second-order accuracy while remaining robust close to singular points where the field becomes stiff. Integration stops as soon as the streamline hits a domain boundary, another singularity, or a c0 corner where the cross-field is discontinuous. The implementation lives in StreamlineGenerator_v2.runge_kutta_heun_integrate_streamline. Kowalski Proposition 9 guarantees the separatrix count: index plus one yields exactly three arms, index minus one yields exactly five arms.
}

\end{frame}

\end{document}
